% GENERATED, edit the French unit or the translation table
\mxnewpage\section*{EV2-T1 -- Mock exam 1}\label{EV2-T1}
\addcontentsline{toc}{section}{EV2-T1 -- Mock exam 1}
\textbf{Practice paper}. Chapters 1 to 7.\par Duration: 2 h. No calculator. Marked out of 20. 14--16 January 2027.
\subsection*{Part 1 -- Fluency (6 points)}
Twelve questions, 0.5 point each. Only one correct answer; no justification required. An incorrect or missing answer scores 0, with no loss of points.
\begin{enumerate}
\item Calculate $2+2\times3$.\par A. $8$\quad B. $12$\quad C. $11$\quad D. $7$
\item Calculate $\frac12+\frac13$.\par A. $\frac{1}{6}$\quad B. $\frac{5}{3}$\quad C. $\frac{2}{5}$\quad D. $\frac{5}{6}$
\item Calculate $\dfrac{2^2}{2^{0}}$.\par A. $2$\quad B. $7$\quad C. $1$\quad D. $4$
\item A value of 50 increases by 20 \%. Final value?\par A. $40$\quad B. $60$\quad C. $70$\quad D. $10$
\item A value of 80 is obtained after a 20 \% decrease. Initial value?\par A. $96$\quad B. $100$\quad C. $64$\quad D. $60$
\item A 25 \% increase then a 20 \% decrease: overall coefficient?\par A. $\frac{4}{5}$\quad B. $\frac{9}{10}$\quad C. $1$\quad D. $\frac{21}{20}$
\item The solution of $2x+2=0$ is\par A. $-4$\quad B. $-1$\quad C. $1$\quad D. $4$
\item The slope between $(0;2)$ and $(2;8)$ is\par A. $6$\quad B. $2$\quad C. $3$\quad D. $-3$
\item The mean of $2, 4, 6$ is\par A. $4$\quad B. $2$\quad C. $12$\quad D. $5$
\item A sector of $90^\circ$ represents what fraction of the total?\par A. $\frac{1}{3}$\quad B. $\frac{3}{4}$\quad C. $\frac{1}{4}$\quad D. $\frac{1}{2}$
\item A probability is $\dfrac{2}{5}$. Its complement is\par A. $\frac{2}{5}$\quad B. $\frac{5}{2}$\quad C. $\frac{1}{5}$\quad D. $\frac{3}{5}$
\item For $f(x)=x^2+2$, $f(-2)$ is\par A. $4$\quad B. $8$\quad C. $6$\quad D. $-2$
\end{enumerate}
\mxnewpage\subsection*{Part 2 -- Independent exercises (14 points)}
Every step must be justified. Exact results are expected unless stated otherwise.
\mxneed{12}
\subsection*{Exercise 1 -- Manufacturing and quality control (5 points)}
Two workshops A and B produce 30 percent and 70 percent respectively of a company's items. Among the items from A, 10 percent are defective; this proportion is 20 percent for B. An item is chosen at random from the total production. The events $A$ and $D$ mean respectively ``comes from workshop A'' and ``is defective''.
\begin{enumerate}
\item Construct a weighted tree whose first level distinguishes the two workshops. \hfill(1 pt)
\item Calculate $P(D)$, then the probability that a defective item comes from A. \hfill(2 pts)
\item Are the events $A$ and $D$ independent? Justify numerically. \hfill(2 pts)
\end{enumerate}
\mxneed{12}
\subsection*{Exercise 2 -- Minimum of a quotient (5 points)}
The function $f$ is defined on the interval $]0;+\infty[$ by $f(x)=x+\dfrac9x$. We study its variation and certain tangents to its curve. The coordinates requested are to be given in exact form.
\begin{enumerate}
\item Calculate and factorise the derivative. \hfill(1 pt)
\item Study the variation and give the minimum. \hfill(2 pts)
\item Determine the points where the tangent has slope $\dfrac{3}{4}$. \hfill(2 pts)
\end{enumerate}
\mxneed{12}
\subsection*{Exercise 3 -- Projection and circle (4 points)}
In an orthonormal frame, consider the point $P(5;2)$ and the line $d$ with equation $x+y-3=0$. Let $H$ be the orthogonal projection of $P$ onto $d$.
\begin{enumerate}
\item Give a nonzero normal vector to $d$. \hfill(1 pt)
\item Determine the coordinates of $H$ and check the conditions of membership and orthogonality. \hfill(1 pt)
\item Give an equation of the circle with centre $P$ passing through $H$. \hfill(2 pts)
\end{enumerate}
